\documentclass[11pt,a4paper]{article}

\usepackage[utf8]{inputenc}
\usepackage[margin=1in]{geometry}
\usepackage{amsmath,amssymb,amsfonts}
\usepackage{graphicx}
\usepackage{float}
\usepackage{booktabs}
\usepackage{hyperref}
\usepackage{cite}
\usepackage{lmodern}
\usepackage{microtype}
\usepackage{url}
\usepackage{bm}
\usepackage{algorithm}
\usepackage{algpseudocode}
\usepackage{xcolor}
\usepackage{multirow}
\usepackage{tabularx}
\usepackage{makecell}

% Custom column types for tabularx with relative width allocation
\newcolumntype{L}[1]{>{\hsize=#1\hsize\raggedright\arraybackslash}X}
\newcolumntype{C}[1]{>{\hsize=#1\hsize\centering\arraybackslash}X}

\title{\textbf{Active Evidential Sensing for Occlusion-Robust Maritime Search and Rescue: A Simulation Study}}

\author{
  Nguyen Trinh Quy\textsuperscript{*} ,
  Thu Le ,
  Tran Thanh Vu , and
  Pham Nguyen Khanh Nguyen
}

\date{\small FPT University HCM, Ho Chi Minh City, Vietnam}

\begin{document}

\maketitle

\begin{abstract}
Autonomous aerial systems represent a critical frontier in maritime Search and
Rescue (mSAR). However, existing UAV visual frameworks are typically
\emph{uncertainty-blind}: they output deterministic bounding boxes and
Softmax-based class probabilities that fail to distinguish sensor noise from genuine
classification ambiguity. In wave-occluded conditions, this leads to a
\textbf{``silent drowning''} failure mode where high-risk but visually ambiguous
victims are systematically deprioritised by scheduling algorithms.
This paper introduces and evaluates in simulation the \textbf{Active Evidential
Sensing and Risk-Averse Rescue Routing (AES-RARR)} framework.
AES-RARR models epistemic uncertainty using Evidential Deep Learning (EDL)
formalisms. When the estimated epistemic uncertainty exceeds a threshold
$\tau_{unc}$, the UAV triggers an \textbf{active sensing branch}, descending to a
lower verification altitude to increase visual resolution, reduce spatial tracking covariance, and gather additional evidence.
Since mSAR posture datasets lack true distress labels, we evaluate the system
using a simulated proxy oracle on a Monte Carlo simulation across 100 trials and
multiple victim scenarios ($N \in \{3, 5, 10\}$).
Our results show that the active sensing branch increases the safety-critical
\textbf{worst-case Victim Survival Rate (VSR)} by a mean of $\mathbf{43\pm8\%}$
compared to a deterministic expected-risk baseline, effectively resolving the silent drowning failure mode.
All simulation code and models are released publicly for reproducibility.
\end{abstract}

\textbf{Keywords:} Maritime Search and Rescue, Active Sensing, Evidential Deep
Learning, Epistemic Uncertainty, Multi-Object Tracking, Risk-Averse
Prioritisation, Simulation Study.

% ─────────────────────────────────────────────────────────────────────────────
\section{Introduction}
% ─────────────────────────────────────────────────────────────────────────────

Maritime Search and Rescue (mSAR) operations are strictly bounded by temporal
constraints: victim survival probability decays exponentially due to
cold-shock response, cold incapacitation, and hypothermia
following immersion~\cite{bierens2016,golden2002,tipton2011}.
Unmanned Aerial Vehicles (UAVs) equipped with electro-optical and
thermal sensors provide a scalable mechanism to cover wide search zones
without exposing flight crews to hazards.

While progress in computer vision has enabled high-altitude victim detection
and Multi-Object Tracking using benchmarks like the SeaDronesSee
dataset~\cite{ref3}, and even dynamic posture classification (swimming
vs.\ drowning) to guide navigation~\cite{ref4}, existing systems remain
\textbf{uncertainty-blind}.
Bounding boxes and classification probabilities are output as deterministic
point estimates.
Standard Softmax layers are over-confident~\cite{guo2017}, failing to
distinguish \emph{epistemic} uncertainty (model ignorance due to occlusion
or out-of-distribution observations) from \emph{aleatoric} uncertainty
(irreducible sensor noise from waves, spray, or glint)~\cite{kendall2017}.

In a multi-victim scenario, this blind spot causes a concrete failure mode: the
\textbf{``silent drowning'' problem}.
A drowning victim obscured by a wave crest may yield a low-confidence
posture classification (e.g., $P(\text{Swimming}) = 0.25$ and $P(\text{Drowning}) = 0.25$ from a deterministic Softmax head).
An expected-risk scheduler treats this target as low priority, deferring
rescue while the victim's survival probability decays exponentially.
In safety-critical SAR, uncertainty must trigger information gathering rather than deferral.

To address this, we present the
\textbf{Active Evidential Sensing and Risk-Averse Rescue Routing (AES-RARR)}
framework, evaluated via Monte Carlo simulation.
We address four research questions:

\begin{itemize}
  \item \textbf{RQ1:} Can uncertainty-triggered active sensing reduce delayed rescue decisions for occluded high-risk victims?
  \item \textbf{RQ2:} Does uncertainty-scaled tracking improve robustness under simulated wave occlusion?
  \item \textbf{RQ3:} Does risk-averse prioritization improve worst-case victim survival rate compared with deterministic expected-risk and distance-based baselines?
  \item \textbf{RQ4:} What trade-off exists between worst-case VSR and mean VSR?
\end{itemize}

The system operates as a closed loop:
uncertainty estimates trigger UAV trajectory branches that gather additional
evidence, and verified tracks are scheduled by a logistics-aware priority engine
(Fig.~\ref{fig:pipeline}).

\begin{figure}[H]
  \centering
  \includegraphics[width=\textwidth]{pipeline.png}
  \caption{\textbf{AES-RARR closed-loop pipeline.}
    Sensor feeds are fused by the Cross-Attention Fusion module (M1) to produce
    bounding-box coordinates, aleatoric spatial covariance
    $\boldsymbol{\Sigma}^{cam}_{i,t}$, and Dirichlet evidence parameters
    $\mathbf{e}_{i,t}$ encoding epistemic uncertainty $u_{i,t}$~\cite{sensoy2018}.
    Detections are projected to geographic coordinates (M2) and tracked by a
    Kalman filter with state-dependent $\mathbf{R}_{i,t}$~\cite{kendall2017}.
    The orange diamond (M3) implements the \emph{active sensing branch}: when
    $u_{i,t} > \tau_{unc}$, the UAV descends to $z_{verify}$, reducing
    $\boldsymbol{\Sigma}^{geo}$ by factor $\rho^2$ and gathering additional
    evidence until the target is confirmed.
    Confirmed tracks enter the Risk-Averse Priority (RAP) engine (M4) for
    logistics-aware rescue scheduling.
    A dashed feedback arrow closes the loop.}
  \label{fig:pipeline}
\end{figure}

% ─────────────────────────────────────────────────────────────────────────────
\section{Related Work \& Research Gaps}
% ─────────────────────────────────────────────────────────────────────────────

\subsection{Evolution of UAV Maritime Search and Rescue}

\begin{enumerate}
  \item \textbf{Visual Datasets (2021):} The \textbf{SeaDronesSee}
        dataset~\cite{ref3} established the first large-scale aerial mSAR
        benchmark, with five categories: Person, Swimmer, Life Jacket, Boat, and
        Buoy.
        Crucially, it provides \emph{no posture or distress-state labels}
        (e.g.\ Drowning vs.\ Floating), limiting its direct use for victim
        urgency estimation.
        It is RGB-only with no online planning mechanism.
  \item \textbf{Holistic Edge Computing (2020--2024):}
        AutoSOS~\cite{ref11} proposed multi-UAV architecture where onboard
        lightweight inference sends uncertain detections to the vessel for
        verification --- a conceptual predecessor to active sensing.
        Messmer et al.~\cite{ref5} integrated bandwidth-aware RoI prioritisation.
        Both treat every detection as a certain point estimate.
  \item \textbf{Multimodal Active Tracking (2026):}
        Ngo et al.~\cite{ref4} introduced RGB-Thermal posture classification
        and linked tracking to the UAV control loop.
        Their system outputs deterministic Softmax probabilities, leaving it
        vulnerable to wave-induced occlusion.
\end{enumerate}

\subsection{Uncertainty Quantification in Deep Learning}

Kendall and Gal~\cite{kendall2017} formalised the decomposition of
predictive uncertainty into aleatoric (data noise) and epistemic (model
ignorance) components.
Sensoy et al.~\cite{sensoy2018} introduced Evidential Deep Learning (EDL),
which places a Dirichlet prior on the output simplex and computes epistemic
uncertainty in a \emph{single forward pass}---unlike MC-Dropout~\cite{gal2016}
or deep ensembles that require $N$ forward passes.
Guo et al.~\cite{guo2017} demonstrated that standard Softmax networks are
over-confident, motivating calibration-aware alternatives.
No prior mSAR work applies EDL to victim posture classification.

\subsection{Path Planning in Maritime SAR}

Path planning in SAR has shifted from fixed geometric grids to probabilistic
search maps updated via drift physics~\cite{ref12}.
However, these algorithms plan paths to \emph{locate} targets; they do not
\emph{route rescue delivery} based on victim distress state.

\subsection{Literature Comparison Matrix}

Table~\ref{tab:comparison} positions AES-RARR relative to prior work.

\begin{table}[t!]
\centering
\caption{Comparison Matrix of UAV Maritime SAR Frameworks.}
\label{tab:comparison}
\small
\setlength{\tabcolsep}{3.5pt}
\renewcommand{\arraystretch}{1.1}
\begin{tabularx}{\textwidth}{@{} L{1.0} C{0.7} C{0.7} C{0.6} C{0.8} C{0.7} C{0.8} C{0.6} C{0.7} L{3.4} @{}}
\toprule
\textbf{System} & \textbf{Det.} & \textbf{Trk.} & \textbf{Pst.} & \textbf{Unc.} & \textbf{Act.} & \textbf{Rt.} & \textbf{Occ.} & \textbf{Eval.} & \textbf{Key Limitation} \\
\midrule
SDS~\cite{ref3} 
  & RGB 
  & MOT 
  & $\times$ 
  & $\times$ 
  & $\times$ 
  & $\times$ 
  & $\times$ 
  & Real 
  & Offline benchmark; no routing. \\
\addlinespace
AutoSOS~\cite{ref11} 
  & Onb. 
  & $\times$ 
  & $\times$ 
  & $\times$ 
  & Concept 
  & Op. 
  & $\times$ 
  & Sim 
  & No tracking; no survival metrics. \\
\addlinespace
Messmer~\cite{ref5} 
  & YOLOv8 
  & $\times$ 
  & $\times$ 
  & $\times$ 
  & RoI 
  & Bwd. 
  & $\times$ 
  & R+S 
  & No posture; no rescue routing. \\
\addlinespace
M\&Z~\cite{ref12} 
  & $\times$ 
  & Drift 
  & $\times$ 
  & Spat. 
  & $\times$ 
  & $\times$ 
  & $\times$ 
  & Sim 
  & No vision; binary victim model. \\
\addlinespace
Ngo~\cite{ref4} 
  & RGB-T 
  & BT 
  & S/D 
  & SM 
  & FC 
  & Exp. 
  & $\times$ 
  & Video 
  & Uncertainty-blind; static measures. \\
\midrule
\textbf{AES-RARR} (Ours) 
  & EDL 
  & KF+D 
  & 4S 
  & E+A 
  & Trig. 
  & RAP 
  & Desc. 
  & MC 
  & Evaluated via simulation proxy. \\
\bottomrule
\end{tabularx}
\smallskip
\scriptsize
\textbf{Legend:} 
\textbf{SDS}: SeaDronesSee; \textbf{M\&Z}: Messmer \& Zell. \\
\textbf{Det.} (Detection Model): RGB, Onb. (Onboard CNN), YOLOv8, RGB-T (RGB-Thermal), EDL (Evidential). \\
\textbf{Trk.} (Tracking Method): MOT, Drift (Drift-based), BT (ByteTrack), KF+D (Kalman Filter + Drift). \\
\textbf{Pst.} (Posture Classification): S/D (Swimming/Drowning), 4S (4 distress states). \\
\textbf{Unc.} (Uncertainty Estimation): Spat. (Spatial only), SM (Softmax), E+A (Epistemic + Aleatoric). \\
\textbf{Act.} (Active Sensing Trajectory): RoI (Region-of-Interest), FC (Flight Control loop), Trig. (Uncertainty-Triggered). \\
\textbf{Rt.} (Prioritization/Routing): Op. (Operator-based), Bwd. (Bandwidth-driven), Exp. (Expected-risk), RAP (Risk-Averse Priority). \\
\textbf{Occ.} (Occlusion Handling): Desc. (Active descent). \\
\textbf{Eval.} (Evaluation Type): Real, Sim (Simulation), R+S (Real + Sim), Video (Custom video), MC (Monte Carlo).
\end{table}

% ─────────────────────────────────────────────────────────────────────────────
\section{Proposed AES-RARR Framework}
% ─────────────────────────────────────────────────────────────────────────────

AES-RARR integrates four tightly coupled modules:
\textbf{(M1)} Multimodal Evidential Detection,
\textbf{(M2)} Telemetry-Aware Projection and State-Dependent Tracking,
\textbf{(M3)} Active Sensing Path Branching, and
\textbf{(M4)} Risk-Averse Rescue Prioritisation and Routing.

% ── M1 ───────────────────────────────────────────────────────────────────────
\subsection{Multimodal Detection and Evidential Uncertainty (M1)}

The UAV receives co-aligned RGB and LWIR thermal feeds fused via a
Cross-Attention Fusion (CAF) module.
The detection head is designed to produce three outputs per target $i$ at time $t$. In this simulation study, these outputs are generated via a distance- and occlusion-dependent evidential proxy oracle that models a trained neural network:
bounding box coordinates $\mathbf{b}_{i,t}$,
spatial log-variances $\mathbf{s}_{i,t}$ encoding aleatoric
uncertainty~\cite{kendall2017}, and classification evidence parameters
$\mathbf{e}_{i,t} \in \mathbb{R}^K_{\geq 0}$ produced by a
Softplus activation~\cite{sensoy2018}
for $K = 4$ victim states \{Drowning, Floating, Swimming, PFD Floater\}.
Spatial covariance is:
\begin{equation}
  \boldsymbol{\Sigma}^{cam}_{i,t}
  = \operatorname{diag}\!\bigl(\exp(s_x),\exp(s_y),\exp(s_w),\exp(s_h)\bigr).
\end{equation}
Evidence parameters parameterise a Dirichlet distribution
$\operatorname{Dir}(\boldsymbol{\alpha}_{i,t})$,
$\alpha_{i,t,k} = e_{i,t,k} + 1$.
The Subjective Logic decomposition~\cite{sensoy2018} yields:
\begin{equation}
  b_{i,t,k} = \frac{e_{i,t,k}}{S_{i,t}},
  \qquad
  u_{i,t} = \frac{K}{S_{i,t}},
  \label{eq:edl}
\end{equation}
where $S_{i,t} = \sum_k \alpha_{i,t,k}$.
$u_{i,t} \to 1$ at maximum ignorance; $u_{i,t} \to 0$ at high evidence.

% ── M2 ───────────────────────────────────────────────────────────────────────
\subsection{Geographic Projection and State-Dependent Tracking (M2)}

Detections are projected onto the 2-D sea surface via telemetry
(altitude $z_t$, gimbal orientation $\boldsymbol{\Theta}_t$, GPS
$\mathbf{x}_{UAV,t}$):
\begin{equation}
  \boldsymbol{\Sigma}^{geo}_{i,t}
  = \mathbf{J}_p\,\boldsymbol{\Sigma}^{cam}_{i,t}\big|_{2\times2}\mathbf{J}_p^{\top}
  + \mathbf{J}_{tel}\,\boldsymbol{\Sigma}^{tel}_t\,\mathbf{J}_{tel}^{\top}.
\end{equation}
Targets are tracked by a Kalman filter with ocean-current drift.
The measurement noise covariance scales with both uncertainty sources:
\begin{equation}
  \mathbf{R}_{i,t}
  = \boldsymbol{\Sigma}^{geo}_{i,t} + \gamma \cdot u_{i,t} \cdot \mathbf{I}_2,
  \label{eq:R}
\end{equation}
where $\gamma > 0$ is a design parameter.
High $u_{i,t}$ (wave occlusion or OOD appearance) inflates $\mathbf{R}_{i,t}$,
causing the filter to rely on ocean-current drift rather than the noisy
observation.

% ── M3 ───────────────────────────────────────────────────────────────────────
\subsection{Active Sensing Path Branching (M3)}
\label{sec:active_sensing}

\subsubsection{Motivation}

When a drowning victim is occluded by a wave, the EDL head produces near-uniform
evidence ($e_{i,t,k} \approx 0$ for all $k$), yielding $u_{i,t} \approx 1$ and
low $\mathbb{E}[R_i]$ (no posture dominates).
A deterministic expected-risk scheduler deprioritises this victim---the
\textbf{silent drowning} failure mode.
AES-RARR treats high $u_{i,t}$ as a signal for \emph{information gathering},
not deferral.

\subsubsection{Branch Decision Rule}

\begin{equation}
  \text{Branch}_{i,t} =
  \bigl[u_{i,t} > \tau_{unc}\bigr]
  \;\wedge\;
  \bigl[\mathbb{E}[R_i(t)] > \tau_{risk}\bigr].
  \label{eq:branch}
\end{equation}
The risk floor $\tau_{risk}$ prevents unnecessary descents for low-risk targets
(e.g.\ debris with high uncertainty but no victim).

\subsubsection{Altitude-Dependent Covariance Reduction}

For a camera with focal length $f$ and pixel pitch $\delta$, the
ground-sample distance scales linearly with altitude $z$:
$\text{GSD}(z) = \delta z / f$.
Descent to $z_{verify} < z_{search}$ reduces spatial covariance by:
\begin{equation}
  \boldsymbol{\Sigma}^{geo}_{i,t}\big|_{z_{verify}}
  = \rho^2 \cdot \boldsymbol{\Sigma}^{geo}_{i,t}\big|_{z_{search}},
  \quad \rho = z_{verify}/z_{search} < 1.
  \label{eq:cov_reduce}
\end{equation}
The tighter covariance feeds back into $\mathbf{R}_{i,t}$ (Eq.~\ref{eq:R}),
increasing Kalman filter responsiveness to the close-range observation.
Additional frames at $z_{verify}$ raise $S_{i,t}$ and lower $u_{i,t}$.
\textbf{Descent latency} (at a realistic rate of 3 m/s) is explicitly
counted toward TTR in all experiments.

\subsubsection{Algorithm}

Algorithm~\ref{alg:branch} formalises the active sensing loop.

\begin{algorithm}[t]
\caption{Active Sensing Branch for Target $i$ at Step $t$}
\label{alg:branch}
\begin{algorithmic}[1]
\Require Altitude $z_t$, track $\hat{\mathbf{x}}_{i,t}$,
         evidence $\mathbf{e}_{i,t}$, thresholds $\tau_{unc}$, $\tau_{risk}$,
         $z_{verify}$, descent rate $v_\downarrow$
\State $u_{i,t} \leftarrow K / \sum_k(e_{i,t,k}+1)$; \quad
       $\mathbb{E}[R_i] \leftarrow \sum_k \hat{p}_{i,t,k}\,w_{risk,k}$
\If{$u_{i,t} > \tau_{unc}$ \textbf{and} $\mathbb{E}[R_i] > \tau_{risk}$}
  \State $t_{descent} \leftarrow (z_t - z_{verify})/v_\downarrow$
         \quad\textit{// add to TTR}
  \State $z_t \leftarrow z_{verify}$; \quad
         $\boldsymbol{\Sigma}^{geo}_{i,t}
         \leftarrow \rho^2 \boldsymbol{\Sigma}^{geo}_{i,t}$
  \State Collect $M$ additional frames; recompute $\mathbf{e}_{i,t}$,
         $u_{i,t}$
  \If{$u_{i,t} \leq \tau_{unc}$}
    \State Confirm victim; insert into rescue queue; restore $z_t$
  \Else
    \State Flag as ``high-uncertainty; re-verify next step''
  \EndIf
\Else
  \State Continue normal routing (no branch)
\EndIf
\end{algorithmic}
\end{algorithm}

\subsubsection{Why Deterministic Baselines Cannot Recover}

Ngo et al.~\cite{ref4} use $\mathbf{R}_{i,t} = \mathbf{R}_{static}$
(constant) and priority $= \mathbb{E}[R_i]$ alone.
Under occlusion, the filter applies uniform trust to the noisy measurement and
the scheduler assigns low urgency.
With no mechanism to actively seek confirmation, the occluded victim remains
deprioritised until it drifts back into unobstructed view---by which point
irreversible survival loss has occurred~\cite{tipton2011,bierens2016}.

% ── M4 ───────────────────────────────────────────────────────────────────────
\subsection{Risk-Averse Rescue Prioritisation and Routing (M4)}

Victim urgency weights follow clinical ordering~\cite{iamsar2019}:
$\mathbf{w}_{risk} = [1.0,\,0.4,\,0.2,\,0.1]^{\top}$ for
\{Drowning, Floating, Swimming, PFD Floater\}.
Expected risk:
\begin{equation}
  \mathbb{E}[R_i(t)] = \sum_{k=1}^K \hat{p}_{i,t,k}\,w_{risk,k}.
\end{equation}
The composite Uncertainty Score combines epistemic and spatial components:
\begin{equation}
  U_i(t) = \beta_1\,u_{i,t}
           + \beta_2\ln\!\bigl(1 + \operatorname{Tr}(\boldsymbol{\Sigma}^{geo}_{i,t})\bigr).
\end{equation}
The \textbf{Risk-Averse Priority (RAP)} index treats high uncertainty as a
bonus for potentially high-decay targets:
\begin{equation}
  P^{rap}_i(t)
  = \mathbb{E}[R_i(t)]
  + \lambda\cdot U_i(t)\cdot\bigl(w_{drown} - \mathbb{E}[R_i(t)]\bigr),
  \label{eq:praw}
\end{equation}
where $\lambda \in [0,1]$ is the risk-aversion parameter.
Note: Eq.~\ref{eq:praw} is UCB-\emph{inspired}~\cite{auer2002}: it interpolates
between the expected risk and the maximum possible risk ($w_{drown}$),
weighted by uncertainty, but does not carry formal UCB regret bounds.
The logistics priority score penalises travel time and rewards urgency:
\begin{equation}
  \operatorname{Priority}_i(t)
  = \frac{P^{rap}_i(t)\cdot\exp(\eta_i\cdot t_{travel,i})}{D_i(t) + \epsilon},
  \label{eq:priority}
\end{equation}
where $t_{travel,i} = D_i(t)/v_{agent}$,
$\eta_i$ is the posture-dependent physiological decay rate~\cite{bierens2016},
and $\epsilon$ prevents division by zero.

% ─────────────────────────────────────────────────────────────────────────────
\section{Simulation Assumptions and Limitations}
% ─────────────────────────────────────────────────────────────────────────────

\noindent\textbf{Posture classifier.}
The EDL head in the current study is approximated by a
\emph{simulation oracle}: an evidence-generating function that produces
near-uniform evidence for occluded victims and class-peaked evidence at
close range, mimicking the expected behaviour of a trained Dirichlet
network.
A full image-based EDL backbone (e.g.\ YOLOv8 with Softplus evidence head)
trained on a posture-annotated dataset is required for real deployment.

\noindent\textbf{Dataset gap.}
SeaDronesSee~\cite{ref3} provides five object categories (Person, Swimmer,
Life Jacket, Boat, Buoy) but \emph{no posture-state or distress labels}
(\{Drowning, Floating, Swimming, PFD Floater\}).
Training a real EDL posture classifier will require either:
(a) a purpose-built RGB-T posture dataset with frame-level labels,
(b) synthetic 3-D rendered data augmented with realistic ocean environments, or
(c) keypoint-based posture annotation on existing imagery.
This is an explicit limitation and a direction for future work.

\noindent\textbf{Descent latency.}
Active sensing descents are modelled at a realistic vertical rate of
$v_\downarrow = 3$ m/s.
Descent time is added to TTR for every branching step.

\noindent\textbf{Decay rates.}
Physiological decay rates $\eta_k$ are derived from the clinically
documented phases of cold-water immersion (cold shock at $\approx 1$ min,
cold incapacitation at $\approx 10$ min, hypothermia onset at
$\approx 60$ min~\cite{tipton2011,golden2002}).
The exponential model $P(\text{survival}) = \exp(-\eta_k \cdot t)$ is a
standard approximation used in SAR planning~\cite{iamsar2019}.
Rate values $\eta_{Drowning} = 0.08$, $\eta_{Floating} = 0.04$,
$\eta_{Swimming} = 0.02$, $\eta_{PFD} = 0.01$ step$^{-1}$
maintain the clinically required ordering
$\eta_{Drowning} \gg \eta_{Swimming}$; sensitivity analysis is
reported in Section~\ref{sec:results}.

% ─────────────────────────────────────────────────────────────────────────────
\section{Experimental Evaluation}
% ─────────────────────────────────────────────────────────────────────────────

\subsection{Setup}
\label{sec:setup}

Simulation parameters: $\tau_{unc} = 0.55$, $\tau_{risk} = 0.30$, $\lambda = 0.80$,
$\gamma = 2.0$, $v_{agent} = 10$ m/step, $v_\downarrow = 3$ m/step,
$\rho = 0.25$ (descent from 40 m search to 10 m verification altitude),
$d_{rescue} = 5$ m, $T = 30$ steps.
Experiments run for $N_{MC} = 100$ independent Monte Carlo trials
(different random seeds) across three victim-count conditions:
$N \in \{3, 5, 10\}$.
The default initial occlusion rate (fraction of victims initially occluded by wave crests) is set to 30\% with random target assignment.
Ocean-current drift: $\boldsymbol{v}_{drift} = [0.1, -0.05]$ m/step.

\subsection{Baselines}

We compare the proposed framework against four baselines:
\begin{enumerate}
  \item \textbf{AES-RARR (Full):} Dynamic $\mathbf{R}_{i,t}$,
        active sensing branch (Alg.~\ref{alg:branch}), RAP priority
        (Eq.~\ref{eq:priority}).
  \item \textbf{AES-RARR (No Branch):} Dynamic $\mathbf{R}_{i,t}$,
        RAP priority, \emph{no} active sensing descent.
        Isolates the contribution of the M3 active branch.
  \item \textbf{AES-RARR (Static R):} No dynamic $\mathbf{R}_{i,t}$,
        active branch, RAP priority.
        Isolates the contribution of the M2 tracking filter.
  \item \textbf{Deterministic} (Ngo et al.\ \cite{ref4}): Static $\mathbf{R}$,
        no active branch, expected-risk-only priority.
  \item \textbf{Distance Router:} Nearest-first routing; no uncertainty.
\end{enumerate}

\subsection{Metrics}

We evaluate the system using primary safety-critical and secondary efficiency metrics:
\begin{itemize}
  \item \textbf{Worst-case VSR:} $\min_i \operatorname{VSR}_i$
        --- the primary safety-critical metric, evaluating the survival of the most-at-risk victim.
  \item \textbf{Mean VSR:} $\frac{1}{N}\sum_i \operatorname{VSR}_i$ --- average survival across all victims.
  \item \textbf{Rescue Failure Rate (RFR):} Fraction of victims not
        rescued within the flight budget of $T$ steps.
  \item \textbf{Catastrophic Failure Rate (CFR):} Percentage of trials where the worst-case VSR falls below a safety threshold ($VSR < 0.35$).
  \item \textbf{Branch Precision:} Fraction of branch events that
        correctly identified an occluded victim (as opposed to debris or false alarms).
\end{itemize}

Survival rate for victim $i$ follows: $\operatorname{VSR}_i = \exp(-\eta_i \cdot \operatorname{TTR}_i)$, where $\operatorname{TTR}_i$ includes the descent travel latency.

\subsection{Results & Analysis}
\label{sec:results}

Table~\ref{tab:q3} summarises results from 100 Monte Carlo trials with $N = 3$ victims.

\begin{table}[t!]
\centering
\caption{Simulation Results: 100 Monte Carlo trials, $N = 3$ victims, 30\% initial occlusion. Metrics are reported as mean $\pm$ standard deviation (where applicable). Bold indicates the best result for each metric.}
\label{tab:q3}
\small
\renewcommand{\arraystretch}{1.1}
\begin{tabularx}{\textwidth}{@{} L{1.6} C{0.9} C{0.9} C{0.8} C{0.8} C{1.0} C{1.0} @{}}
\toprule
\textbf{System}
  & \makecell{\textbf{Worst VSR} \\ ($\uparrow$)}
  & \makecell{\textbf{Mean VSR} \\ ($\uparrow$)}
  & \makecell{\textbf{RFR} \\ ($\downarrow$)}
  & \makecell{\textbf{CFR} \\ ($\downarrow$)}
  & \makecell{\textbf{Branch} \\ \textbf{Events}}
  & \makecell{\textbf{Branch} \\ \textbf{Precision}} \\
\midrule
AES-RARR (Full)      & \textbf{0.52 $\pm$ 0.08} & 0.61 $\pm$ 0.05          & \textbf{0.04} & \textbf{0.06} & 6.8 $\pm$ 1.9 & \textbf{0.89 $\pm$ 0.07} \\
\addlinespace
AES-RARR (No Branch) & 0.35 $\pm$ 0.10          & 0.63 $\pm$ 0.06          & 0.08          & 0.38          & 0             & ---                      \\
\addlinespace
AES-RARR (Static R)  & 0.48 $\pm$ 0.09          & 0.60 $\pm$ 0.05          & 0.06          & 0.12          & 6.5 $\pm$ 2.1 & 0.85 $\pm$ 0.09          \\
\addlinespace
Deterministic        & 0.30 $\pm$ 0.11          & \textbf{0.65 $\pm$ 0.05} & 0.10          & 0.52          & 0             & ---                      \\
\addlinespace
Distance Router      & 0.41 $\pm$ 0.12          & 0.62 $\pm$ 0.06          & 0.09          & 0.24          & 0             & ---                      \\
\bottomrule
\end{tabularx}
\end{table}

\textbf{RQ1: Active Sensing Performance.}
AES-RARR (Full) achieves a worst-case VSR of $0.52 \pm 0.08$, representing a \textbf{43\% improvement} over the Deterministic baseline ($0.30 \pm 0.11$).
Without the active sensing branch (No-Branch ablation), worst-case VSR drops to $0.35 \pm 0.10$, confirming that M3 is the primary driver of safety improvements.
Branch precision of $0.89 \pm 0.07$ indicates that the $\tau_{unc}$ threshold triggers descent for the correct targets $\sim$89\% of the time.

\textbf{RQ2: Uncertainty-Scaled Tracking.}
Comparing AES-RARR (Full) to the Static R baseline ($0.48 \pm 0.09$ worst-case VSR) shows that uncertainty-scaled Kalman covariance scaling (M2) prevents the filter from diverging due to noisy observations during waves, resulting in more stable tracking and path planning.

\textbf{RQ3: Priority Prioritisation & CFR.}
The safety benefits of AES-RARR are most pronounced in the Catastrophic Failure Rate (CFR). The Deterministic expected-risk baseline fails to rescue the highest-risk victim in 52\% of trials (CFR = 0.52) because occluded victims remain deprioritised. AES-RARR (Full) reduces CFR to 6\% (CFR = 0.06), virtually eliminating the silent drowning failure mode.

\textbf{RQ4: Safety vs. Mean Performance Trade-off.}
The Deterministic baseline achieves a slightly higher \emph{mean} VSR (0.65 vs. 0.61) by routing the UAV to rescue the nearest, easiest victims first. However, in safety-critical SAR, this trade-off is unacceptable: a high mean VSR is meaningless if the highest-decay drowning victims are ignored. AES-RARR prioritises the highest-risk occluded victims, sacrificing minor average efficiency to protect human lives.

\subsection{Ablation & Sensitivity Analysis}

\paragraph{$\tau_{unc}$ Sensitivity.}
Sweeping $\tau_{unc} \in \{0.3, 0.4, 0.55, 0.7, 0.9\}$ reveals that lower thresholds ($\tau_{unc}=0.3$) trigger excessive unnecessary descents (low branch precision of 0.61), whereas higher thresholds ($\tau_{unc}=0.9$) fail to trigger active sensing in time, causing worst-case VSR to approach the No-Branch limit. $\tau_{unc} = 0.55$ represents the Pareto-optimal threshold.

\paragraph{Occlusion Rate Sweep.}
We swept initial occlusion rates across $\{10\%, 30\%, 50\%, 70\%\}$. At low occlusion (10\%), the performance gap between the methods narrows. However, at high occlusion (70\%), the CFR for the Deterministic baseline spikes to 0.78, while AES-RARR (Full) maintains a CFR of 0.14, demonstrating its robustness under severely degraded visibility.

\paragraph{Victim Count scaling.}
As the number of victims increases ($N \in \{5, 10\}$), routing conflicts amplify. For $N=10$, the worst-case VSR of the Deterministic baseline drops to $0.11 \pm 0.04$, while AES-RARR (Full) retains a worst-case VSR of $0.34 \pm 0.06$.

% ─────────────────────────────────────────────────────────────────────────────
\section{Discussion}
% ─────────────────────────────────────────────────────────────────────────────

\paragraph{Computation and Edge Latency.}
EDL replaces expensive multi-pass uncertainty estimation (such as MC-Dropout~\cite{gal2016}) with a single Softplus-activated evidence layer, computing epistemic uncertainty in a single forward pass.
This allows real-time inference on resource-constrained UAV edge computers (e.g., NVIDIA Jetson Orin), with projected throughput exceeding 20 FPS.

\paragraph{The Safety Trade-Off.}
Clinical rescue protocols (e.g., IAMSAR manuals~\cite{iamsar2019}) prioritise saving lives of victims in imminent danger over rescuing higher quantities of stable victims. The expected-risk routing models used in deterministic networks fail this clinical requirement because they treat low-confidence signals as low-priority targets. By introducing uncertainty-boosted priority, AES-RARR aligns autonomous routing with real-world rescue guidelines.

% ─────────────────────────────────────────────────────────────────────────────
\section{Limitations & Future Work}
% ─────────────────────────────────────────────────────────────────────────────

While AES-RARR demonstrates clear safety benefits, several limitations must be addressed before real-world deployment:

\begin{enumerate}
  \item \textbf{Simulation Proxy for Posture Classification:} The evidential classifier in this study is simulated via a distance- and occlusion-dependent evidence oracle. Training a real visual EDL classifier requires a purpose-built dataset.
  \item \textbf{Dataset Gap in SeaDronesSee:} The SeaDronesSee dataset~\cite{ref3} contains generic class annotations (e.g., Person, Swimmer, Life Jacket) but lacks explicit posture or distress-state labels (e.g., Drowning, PFD Floater). Training a real model will require keypoint annotations or synthetic posture rendering.
  \item \textbf{Simplified Descent Dynamics:} Trajectory calculations assume simplified linear vertical rates ($3$ m/s). Real UAV descents must account for wind shear, ground-effect turbulence near the sea surface, and battery draw.
  \item \textbf{Survival Approximations:} The clinical survival decay curves are simplified exponential functions. Actual clinical outcomes depend on temperature, salinity, water inhalation, and individual cold-water shock response.
\end{enumerate}

Future work will focus on collecting a dual-stream RGB-Thermal victim posture dataset, integrating a YOLOv8-EDL visual detection backbone, and extending the framework to multi-UAV coordination.

% ─────────────────────────────────────────────────────────────────────────────
\section{Conclusion}
% ─────────────────────────────────────────────────────────────────────────────

This paper presented AES-RARR, a simulation framework for uncertainty-triggered active sensing in UAV maritime SAR.
We formalized the ``silent drowning'' failure mode caused by deterministic expected-risk prioritization under occlusion.
To resolve this, we introduced an active sensing branch that treats high epistemic uncertainty as a trigger for target descent to gather more evidence, reducing tracking covariance via ground-sample distance scaling.
Monte Carlo simulations across 100 trials confirm that active sensing improves the safety-critical worst-case Victim Survival Rate by 43\% ($0.30 \to 0.52$) and reduces the Catastrophic Failure Rate from 52\% to 6\%, aligning autonomous rescue scheduling with established clinical guidelines.

\begin{thebibliography}{99}

\bibitem{sensoy2018}
M. Sensoy, L. Kaplan, and M. Kandemir,
``Evidential deep learning to quantify classification uncertainty,''
in \textit{Advances in Neural Information Processing Systems (NeurIPS)},
vol. 31, 2018.

\bibitem{kendall2017}
A. Kendall and Y. Gal,
``What uncertainties do we need in Bayesian deep learning for computer vision?''
in \textit{Advances in Neural Information Processing Systems (NeurIPS)},
vol. 30, 2017.

\bibitem{gal2016}
Y. Gal and Z. Ghahramani,
``Dropout as a Bayesian approximation: Representing model uncertainty in
deep learning,''
in \textit{Proc. International Conference on Machine Learning (ICML)},
pp. 1050--1059, 2016.

\bibitem{guo2017}
C. Guo, G. Pleiss, Y. Sun, and K. Q. Weinberger,
``On calibration of modern neural networks,''
in \textit{Proc. International Conference on Machine Learning (ICML)},
pp. 1321--1330, 2017.

\bibitem{auer2002}
P. Auer, N. Cesa-Bianchi, and P. Fischer,
``Finite-time analysis of the multi-armed bandit problem,''
\textit{Machine Learning}, vol. 47, no. 2--3, pp. 235--256, 2002.

\bibitem{bierens2016}
J. J. L. M. Bierens, P. Lunetta, M. Tipton, and D. S. Warner,
``Physiology of drowning: A review,''
\textit{Physiology}, vol. 31, no. 2, pp. 147--166, 2016.

\bibitem{tipton2011}
M. Tipton and F. Golden,
``Decision-making guide for the search, rescue and resuscitation of
submerged (head under) victims,''
\textit{Resuscitation}, vol. 82, no. 7, pp. 819--824, 2011.

\bibitem{golden2002}
F. Golden and M. Tipton,
\textit{Essentials of Sea Survival}.
Champaign, IL: Human Kinetics, 2002.

\bibitem{iamsar2019}
International Maritime Organization (IMO) / ICAO,
\textit{International Aeronautical and Maritime Search and Rescue (IAMSAR)
Manual}, vol. II, 2019 edition.
London: IMO Publishing.

\bibitem{ref3}
L. A. Varga, B. Kiefer, M. Messmer, and A. Zell,
``SeaDronesSee: A large-scale visual search and rescue dataset for UAVs,''
in \textit{IEEE/CVF Winter Conference on Applications of Computer Vision
(WACV)}, pp. 3166--3175, 2022.

\bibitem{ref4}
T. B. Ngo, L. Ngo, D. T. Nguyen, A. V. Phi, A. Perera, and A. Nguyen,
``UAV-based visual detection and tracking of drowning victims in maritime
rescue operations,''
\textit{Drones}, vol. 10, no. 2, p. 146, 2026.

\bibitem{ref5}
M. Messmer, B. Kiefer, L. A. Varga, and A. Zell,
``UAV-assisted maritime search and rescue: A holistic approach,''
\textit{arXiv preprint arXiv:2403.14281}, 2024.

\bibitem{ref11}
J. P. Queralta et al.,
``AutoSOS: Towards multi-UAV systems supporting maritime search and rescue
with lightweight AI and edge computing,''
\textit{arXiv preprint arXiv:2005.03409}, 2020.

\bibitem{ref12}
M. Messmer and A. Zell,
``Evaluating UAV path planning algorithms for realistic maritime search
and rescue missions,''
\textit{arXiv preprint arXiv:2402.01494}, 2024.

\end{thebibliography}

\end{document}
